\section{LAN8660/LAN8661 Development Platform}
\label{sec:lan866x}

\begin{risknote}[title={Verified against primary Microchip sources -- the named parts are not what the block diagrams assume}]
Microchip's LAN8660/LAN8661/LAN8662 are the \textbf{VelocityDRIVE 10BASE-T1S Endpoint family} (datasheet/sell-sheet DS00006257C, current as of this planning pass). They are explicitly a \emph{``software-less architecture''}: a 32-pin (5$\times$5\,mm) VQFN part combining a 10BASE-T1S PMD transceiver with a fixed-function ``EP'' (endpoint) block that exposes I$^2$C/SPI/UART pins directly to peripherals, configured, controlled, and diagnosed entirely over the network via the standards-based Remote Control Protocol (RCP) and Microchip's MPLAB Network Creator GUI tool -- \emph{not} by a customer-written embedded application. \textbf{LAN8660} is the ``Control endpoint'' (sensors/actuators); \textbf{LAN8661} is the ``Lighting endpoint'' (bridges T1S video/control data to LED driver ICs); \textbf{LAN8662} (Audio endpoint, mic/speaker bridge) is a sibling part not named in the original brief. All three are AEC-Q100 Grade~1 ($-40$ to $+125^\circ$C) and ISO~26262 ASIL~B qualified.

This does not match the common-node block diagram in Section~\ref{sec:common-platform} (MCU $\rightarrow$ SPI $\rightarrow$ T1S network chip $\rightarrow$ 10BASE-T1S), because that diagram assumes a general-purpose network component sitting \emph{behind} a customer MCU running Alfred's own firmware. The part family that actually matches that architecture is \textbf{Microchip LAN8650/LAN8651}: a MAC-PHY (MAC~+~PHY combined) 10BASE-T1S device with a standard SPI host interface implementing the OPEN Alliance TC6 10BASE-T1x MAC-PHY Serial Interface specification (SPI clock up to 25\,MHz), designed specifically so ``low-cost microcontrollers, including those without an onboard MAC,'' can reach a 10BASE-T1S network. LAN8651 adds an integrated LDO for 3.3\,V-only designs; LAN8650 requires external supplies. Both support PLCA (node~0 = coordinator, followers numbered 1--0xFE, transmit-opportunity count configurable, practically $\geq$8 followers).

\textbf{Implication for this program:} the common node's shared network block (Section~\ref{sec:common-platform}) should be built around \textbf{LAN8650/LAN8651}, not LAN8660/LAN8661 -- the Gateway role structurally requires a host MCU running Alfred's own capture/timestamp/normalize/mirror/translate logic (Section~\ref{sec:arch-b}), and LAN8660/8661 cannot host that logic or speak CAN/CAN-FD/LIN at all. LAN8660/8661 remain in scope as a separate, explicitly time-boxed investigation (Section~\ref{sec:rev-a}) into a possible \emph{second, MCU-less} lightweight E2B product for the simplest peripheral cases -- not as the basis for Rev~A/B/C.
\end{risknote}

\subsection{Family comparison}

\begin{center}
\small
\begin{tabularx}{\linewidth}{@{}p{2.4cm} p{2.6cm} X@{}}
\toprule
\textbf{Part} & \textbf{Class} & \textbf{Role in this program} \\
\midrule
LAN8650 / LAN8651 & MAC-PHY, SPI host interface (OA TC6), PLCA & \textbf{Primary.} Network component for the common node in both E2B and Gateway configurations. LAN8651 preferred for prototyping (integrated LDO simplifies the 3.3\,V rail). \\
LAN8660 & VelocityDRIVE control endpoint, MCU-less & \textbf{Secondary, time-boxed.} Candidate for a future MCU-less lightweight E2B SKU; characterize on Rev~A only if schedule allows (Section~\ref{sec:rev-a}). \\
LAN8661 & VelocityDRIVE lighting endpoint, MCU-less & \textbf{Out of scope for Sept--Nov.} Lighting-specific bridge to LED driver ICs; no CAN/LIN path, no generic peripheral I/O -- does not serve either target architecture directly. \\
LAN8670 / LAN8671 / LAN8672 & PHY-only (RMII/MII to a host MAC) & \textbf{Reference/comparison only.} Useful as a known-good third-party T1S node (via Microchip's own EVB-LAN8670-USB) for bus-side validation, not as the common node's network component, since it needs a host with its own Ethernet MAC rather than a simple SPI link. \\
\bottomrule
\end{tabularx}
\end{center}

\subsection{Design and layout notes (LAN8650/LAN8651)}

Host interface is 4-wire SPI (OA TC6 chunk/footer protocol) plus an interrupt pin, a general reset, and (on the MikroE reference layout) GPIO on DA0/WK0 -- fits comfortably on the common node's MCU SPI peripheral with one interrupt line and one GPIO reset. PLCA configuration (node ID, transmit-opportunity count) is a register write at init, per Microchip's LAN8650/1 Configuration Application Note (DS60001760) -- no special layout impact beyond standard SPI routing discipline. T1S coupling follows the standard transformer-coupled reference design; the two-position termination jumper pattern seen on Microchip/MikroE reference boards (both populated only at physical segment ends) is the layout pattern to copy for Rev~A/B. Clock: onboard crystal per the MAC-PHY's own reference design, no exotic requirements. Errata/availability: no blocking errata identified against current (September 2026) documentation; part is in production and stocked by MikroE-affiliated and standard distribution.

\subsection{Evaluation hardware to purchase immediately}

\begin{center}
\small
\begin{tabularx}{\linewidth}{@{}p{3.6cm} X@{}}
\toprule
\textbf{Board} & \textbf{Why} \\
\midrule
MikroE Two-Wire ETH Click (MIKROE-5543) $\times$3 & LAN8651B1 MAC-PHY, SPI + IRQ + GPIO + reset on a mikroBUS footprint -- plugs \emph{directly} into the Clicker~4's mikroBUS sockets with zero custom hardware. This is the fastest path to ``bring up first T1S development link'' in Week~1 (Section~\ref{sec:schedule}). Buy at least 3 so a 2--3 node multidrop segment can be assembled from click boards alone before Rev~A silicon ever arrives. \\
MikroE Two-Wire ETH 3~Click (MIKROE-6550) & Newer-generation LAN8651B1 click board; buy 1--2 as a hedge against MIKROE-5543 stock/revision differences and to cross-check PLCA behavior between board revisions. \\
Microchip EVB-LAN8670-USB & PHY-only family, USB-attached -- useful as an independent, known-good bus-side reference node/observation point that does not depend on Alfred's own firmware, for isolating ``is this a network problem or a firmware problem'' during Rev~A/B bring-up. \\
LAN8650/LAN8651 samples (loose die-form parts, several units) & For Rev~A/B schematic capture and BOM, ordered directly rather than only relying on click boards for the custom PCB work. \\
\bottomrule
\end{tabularx}
\end{center}

Software support: Microchip publishes an open OA-TC6 driver (\texttt{MicrochipTech/oa-tc6-lib} on GitHub) with example applications; upstream driver support also exists in Zephyr and in Espressif's component registry, both useful as reference implementations even though Alfred's own MCU choice (Section~\ref{sec:rev-b}) may differ from either ecosystem. MPLAB Network Creator (the GUI tool for the VelocityDRIVE endpoint family) is relevant only to the secondary LAN8660 investigation, not to the LAN8650/1-based common node firmware.
